\par\SourceLeaf{528}Man kann die Vorschrift also auch so ausdrücken; (immer voraussetzend $p\equiv5\ (\mathrm{mod}.8)$)
\[
\begin{aligned}
\text{Es ist }b&\equiv a+1\ (\mathrm{mod}.8),&\quad&\text{wenn }m\text{ gerade}\\
b&\equiv a+5&&\text{wenn }m\text{ ungerade.}
\end{aligned}
\]
Ich wage noch keine Vermuthung, ob ein noch einfacheres Criterium möglich ist, woran man den Fall des geraden $m$ von dem des ungeraden \textit{im Voraus} unterscheiden könnte, d. i. ohne den Werth von $m$ selbst zu kennen, da, wie ich schon oben bemerkt habe, dies Rapprochement noch ganz neu ist.
\par Für den Fall $p\equiv1\ (\mathrm{mod}.8)$, bleibt zwar obige Congruenz $b\equiv2k+a-1\ (\mathrm{mod}.8)$ richtig, entscheidet aber nicht mehr über das Zeichen von $b$, da sie dem positiven und negativen Werthe von $b$ zugleich genug thut. Es ist hier nemlich $k$ immer gerade, $=2m$ (wenn die Bedeutung von $m$ eben so ausgesprochen wird wie oben) oder $=2m+2$, wenn man unter $m$ die Anzahl der Zerlegungen von $p$ in 3 positive \textit{ungleiche} Quadrate versteht, und $b\equiv0\ (\mathrm{mod}.4)$, oder $b\equiv-b\ (\mathrm{mod}.8)$. Ich vermuthe dass der Fall $p\equiv1\ (\mathrm{mod}.8)$ oder $b\equiv0\ (\mathrm{mod}.4)$ altioris indaginis ist und vielleicht wieder
\[
\begin{aligned}
b&\equiv4\ (\mathrm{mod}.8)&&\text{leichter als }b\equiv0\ (\mathrm{mod}.8)\\
b&\equiv8\ (\mathrm{mod}.16)&&\text{leichter als }b\equiv0\ (\mathrm{mod}.16)
\end{aligned}
\]
u. s. w.
\par\Needspace{5\baselineskip}Mit ausgezeichneter Hochachtung beharre ich
\par\nobreak\noindent\hfill Ihr freundschaftlich ergebenster
\par\nobreak\noindent Göttingen den 30. Mai 1828.\hfill C. F. \textsc{Gauss}.
\BookRule{.25\textwidth}
